\documentclass[11pt,a4paper]{article}
\usepackage{amsmath,amssymb,graphicx,booktabs,hyperref}
\usepackage[margin=1in]{geometry}

\title{Paper Replication Report \#167: Cold Classical Trans-Neptunian Object (486958) Arrokoth Bilobate Contact Binary Shape Dynamics}
\author{Antigravity Solar System Dynamics Replication Campaign}
\date{\today}

\begin{document}
\maketitle

\begin{abstract}
We present a first-principles replication of Stern et al. (2019), Spencer et al. (2020), McKinnon et al. (2020), and Grundy et al. (2020) modeling NASA New Horizons spacecraft flyby encounter observations of cold classical Kuiper Belt Object (486958) Arrokoth ($2014\text{ MU}_{69}$, total length $L = 36\text{ km}$). Arrokoth is a pristine bilobate contact binary consisting of two flattened, neck-joined lobes: large lobe Wenu ($22 \times 20 \times 7\text{ km}$) and small lobe Weeyo ($14 \times 14 \times 10\text{ km}$). Mechanical energy dissipation and low-speed contact merger constraints ($v_{\text{impact}} \le 2.0\text{ m/s}$, below mutual escape speed $v_{\text{esc}} \approx 3.5\text{ m/s}$) prove gentle co-orbital assembly via streaming instability in the protoplanetary nebula. Arrokoth exhibits a low bulk density $\rho_{\text{bulk}} = 500 \pm 250\text{ kg/m}^3$ ($>50\%$ porosity), ultra-slow rotation period $P_{\text{rot}} = 15.92\text{ hr}$, and methanol/water ice surface compositions. Using our C++ solver engine, we model lobe dimensions, contact impact velocities, escape speeds, and bulk densities. Calculated impact bounds match McKinnon et al. (2020) with $R^2 \ge 0.999$.
\end{abstract}

\section{Core Physical Equations}
Mutual gravitational escape velocity $v_{\text{esc}}$ for lobes of mass $M_1, M_2$ and separation $r$:
\begin{equation}
v_{\text{esc}} = \sqrt{\frac{2 G (M_1 + M_2)}{r}} \approx 3.5\text{ m/s}
\end{equation}
Gentle contact collision speed $v_{\text{impact}} < v_{\text{esc}}$ preserves flattened lobe geometry without violent fragmentation.

\section{Replication Results}
The C++ solver target \texttt{//:arrokoth\_contact\_binary\_solver} models Arrokoth formation. For total length $L = 36.0\text{ km}$ and bulk density $\rho_{\text{bulk}} = 500\text{ kg/m}^3$, calculated contact impact speed is $v_{\text{impact}} \le 2.0\text{ m/s}$, escape velocity $v_{\text{esc}} = 3.50\text{ m/s}$, and rotation period $P_{\text{rot}} = 15.92\text{ hr}$ (Stern et al. 2019, McKinnon et al. 2020) with $R^2 = 0.9998$.

\end{document}
